\SourceLeaf{369}%
\noindent die Cosinus der drei Winkel sein, welche eine Gerade $0P$ resp. mit $01,02,03$ macht. Diese gerade Linie hat in dem bewegten Raume nach der Mutation dieselbe Lage wie vor derselben; [wird ferner]
\[
\BareRoot{(aa+\rho\rho)}=k
\]
gesetzt, oder
\[
k=\BareRoot{(aa+bb+cc+dd)},
\]
ist $k$ die linearische Vergrösserung des mutirten Raumes.

Endlich
\[
a=k\cos\theta,\qquad \rho=k\sin\theta
\]
gesetzt, ist $2\theta$ diejenige Drehung um $0P$, durch welche, verbunden mit der Vergrösserung $k$, die Mutation erzeugt werden kann.

\Needspace{12\baselineskip}
\begin{center}\bfseries BEMERKUNG.\end{center}

Die unter [I] zusammengefassten Notizen stehen in einem Handbuche. Die Nummern 1 bis 7 stimmen in Schrift und Tinte genau mit astronomischen Aufzeichnungen aus dem Jahre 1819 überein, die sich auf den vorhergehenden Seiten befinden. Verschieden von ihnen sieht die als Nummer 8 bezeichnete Bemerkung aus. Sie trägt sachlich den Character eines spätern Zusatzes, und diese Annahme wird auch dadurch bestätigt, dass augenscheinlich gleichzeitig mit ihr niedergeschriebene astronomische Notizen aus der Zeit von 1822 bis 1823 stammen. Die Notiz [II] füllt einen einzelnen Zettel ohne Datirung. Nach der Handschrift zu urtheilen, stammt sie aus noch späterer Zeit.
\begin{flushright}\textsc{Stäckel.}\end{flushright}
